\scanpage{47}
de) l'orientation de $\vec\Delta$, soit un syst. \deleted{\unsure{for}}
d'arêtes \uline{appartenant à l'orientation} du \deleted{triangle}
% slanted
\margin{i.e. (\unsure{intuitivement|orientations}) qu'ils bordent un $=$ triangle\ldots}
$\omega_{\{i,j,\Delta\}}$ de l'arrangem\textsuperscript{t} des trois ps. droites
$\Delta,D,D'$ (ce qui caractérise l'orientation de
$\vec D'$, parmi les deux de $D'$, en termes de
celles de $\Delta,D$).

(A2) Cette prescription — qui définit, \deleted{\unsure{une}} pour
$s,\vec\Delta$ fixés, une permutation $\vec u_{s,\vec\Delta}$ de $\vec I_s$
au dessus de la permutation $u_{s,\vec\Delta}$ de $I_s$)
est bien une permutation \uline{circulaire}
\deleted{\unsure{involution sans être}} i.e. $(\vec u_{s,\vec\Delta})^{\unsure{\nu_s|v_s}}$ (qui est, \deleted{soit}
a priori, soit l'identité, soit l'\deleted{autre} transposition
sur chaque fibre de (26)) \deleted{n'est} est la transposition.

\deleted{(A3) Si $\vec{\Delta'}$ est \unsure{$\to$}}
\uline{NB} On aura \deleted{\unsure{$\vec u$}}
\[\begin{aligned}\vec u_{s,-\vec\Delta}&=\vec u_{s,\vec\Delta}^{-1},\qquad\text{posons}\\
\vec{\unsure{Z|\mathcal Z|\mathfrak Z}}_{s,\Delta}&=\{\vec u_{s,\vec\Delta},\vec u_{s,-\vec\Delta}\}\subset{\unsure{\mathfrak S|\mathscr S|\mathcal C}}_{\vec I_s}
\end{aligned}\tag*{(30)}\]

(A3) Soit $\Delta'$ une autre ps-droite ne passant pas
par $s$, alors on a
\[\vec Z_{s,\Delta'}=\vec Z_{s,\Delta}\qquad(\text{soit }\vec Z_s)\]

Bien entendu, les conditions (A1) à (A3)
sont satisfaites pour un arrangem\textsuperscript{t} \inserted{1-dim.} orienté
de ps. droites \uline{provenant} d'un arrangem\textsuperscript{t} 2-dimensionnel,
ce que permet donc de
% archive sheet number: 46
